\chapter{Unit 1: Fundamentals of D.C. Circuits}

\section{Section A: Multiple Choice Question Practice (60 Curated Problems)}

\noindent\textbf{Q1. The SI unit of electrical resistance is the:} \hfill [Topic: resistance | Easy]
\begin{enumerate}[label=(\alph*),topsep=2pt,itemsep=1pt,partopsep=0pt,parsep=0pt]
    \item {Ohm ( $\Omega$ )}
    \item {Farad ( F )}
    \item {Ampere ( A )}
    \item {Volt ( V )}
\end{enumerate}
\vspace{2mm}

\noindent\textbf{Q2. The resistance of a conductor is directly proportional to its:} \hfill [Topic: resistance | Easy]
\begin{enumerate}[label=(\alph*),topsep=2pt,itemsep=1pt,partopsep=0pt,parsep=0pt]
    \item {Current through it}
    \item {Length}
    \item {Temperature coefficient only}
    \item {Cross-sectional area}
\end{enumerate}
\vspace{2mm}

\noindent\textbf{Q3. According to Ohm's law, the current through a resistor is given by:} \hfill [Topic: ohm's law | Easy]
\begin{enumerate}[label=(\alph*),topsep=2pt,itemsep=1pt,partopsep=0pt,parsep=0pt]
    \item {I = V + R}
    \item {I = V $\cdot$ R}
    \item {I = V R}
    \item {I = R V}
\end{enumerate}
\vspace{2mm}

\noindent\textbf{Q4. A 10 $\Omega$ resistor carries a current of 2 A . The voltage across it is:} \hfill [Topic: ohm's law | Easy]
\begin{enumerate}[label=(\alph*),topsep=2pt,itemsep=1pt,partopsep=0pt,parsep=0pt]
    \item {20 V}
    \item {0.2 V}
    \item {12 V}
    \item {5 V}
\end{enumerate}
\vspace{2mm}

\noindent\textbf{Q5. Voltage is best described as the:} \hfill [Topic: voltage | Easy]
\begin{enumerate}[label=(\alph*),topsep=2pt,itemsep=1pt,partopsep=0pt,parsep=0pt]
    \item {Opposition to current flow}
    \item {Stored electric charge}
    \item {Rate of flow of charge}
    \item {Potential difference between two points}
\end{enumerate}
\vspace{2mm}

\noindent\textbf{Q6. Electric current is defined as the rate of flow of:} \hfill [Topic: current | Easy]
\begin{enumerate}[label=(\alph*),topsep=2pt,itemsep=1pt,partopsep=0pt,parsep=0pt]
    \item {Power}
    \item {Voltage}
    \item {Energy}
    \item {Charge}
\end{enumerate}
\vspace{2mm}

\noindent\textbf{Q7. The SI unit of capacitance is the:} \hfill [Topic: capacitance | Easy]
\begin{enumerate}[label=(\alph*),topsep=2pt,itemsep=1pt,partopsep=0pt,parsep=0pt]
    \item {Ohm ( $\Omega$ )}
    \item {Coulomb ( C )}
    \item {Henry ( H )}
    \item {Farad ( F )}
\end{enumerate}
\vspace{2mm}

\noindent\textbf{Q8. The charge stored on a capacitor is given by:} \hfill [Topic: capacitance | Easy]
\begin{enumerate}[label=(\alph*),topsep=2pt,itemsep=1pt,partopsep=0pt,parsep=0pt]
    \item {Q = C + V}
    \item {Q = C V}
    \item {Q = V C}
    \item {Q = C V}
\end{enumerate}
\vspace{2mm}

\noindent\textbf{Q9. The SI unit of inductance is the:} \hfill [Topic: inductance | Easy]
\begin{enumerate}[label=(\alph*),topsep=2pt,itemsep=1pt,partopsep=0pt,parsep=0pt]
    \item {Weber ( Wb )}
    \item {Farad ( F )}
    \item {Tesla ( T )}
    \item {Henry ( H )}
\end{enumerate}
\vspace{2mm}

\noindent\textbf{Q10. Electrical power dissipated in a resistor can be expressed as:} \hfill [Topic: power and energy concepts | Easy]
\begin{enumerate}[label=(\alph*),topsep=2pt,itemsep=1pt,partopsep=0pt,parsep=0pt]
    \item {P = V I}
    \item {P = V I}
    \item {P = V + I}
    \item {P = I V}
\end{enumerate}
\vspace{2mm}

\noindent\textbf{Q11. The unit of electrical energy is the:} \hfill [Topic: power and energy concepts | Easy]
\begin{enumerate}[label=(\alph*),topsep=2pt,itemsep=1pt,partopsep=0pt,parsep=0pt]
    \item {Ampere ( A )}
    \item {Joule ( J )}
    \item {Volt ( V )}
    \item {Watt ( W )}
\end{enumerate}
\vspace{2mm}

\noindent\textbf{Q12. Kirchhoff's Current Law (KCL) is based on the conservation of:} \hfill [Topic: kirchhoff's laws | Easy]
\begin{enumerate}[label=(\alph*),topsep=2pt,itemsep=1pt,partopsep=0pt,parsep=0pt]
    \item {Energy}
    \item {Voltage}
    \item {Power}
    \item {Charge}
\end{enumerate}
\vspace{2mm}

\noindent\textbf{Q13. Kirchhoff's Voltage Law (KVL) states that the algebraic sum of voltages around a closed loop is:} \hfill [Topic: kirchhoff's laws | Easy]
\begin{enumerate}[label=(\alph*),topsep=2pt,itemsep=1pt,partopsep=0pt,parsep=0pt]
    \item {Equal to the total current}
    \item {Zero}
    \item {Infinite}
    \item {Equal to the source voltage}
\end{enumerate}
\vspace{2mm}

\noindent\textbf{Q14. Three resistors of 2 $\Omega$ , 3 $\Omega$ , and 5 $\Omega$ are connected in series. The total resistance is:} \hfill [Topic: intuitive method of circuit analysis- series and parallel simplification | Easy]
\begin{enumerate}[label=(\alph*),topsep=2pt,itemsep=1pt,partopsep=0pt,parsep=0pt]
    \item {10 $\Omega$}
    \item {30 $\Omega$}
    \item {3.33 $\Omega$}
    \item {0.97 $\Omega$}
\end{enumerate}
\vspace{2mm}

\noindent\textbf{Q15. Two 6 $\Omega$ resistors connected in parallel give a total resistance of:} \hfill [Topic: intuitive method of circuit analysis- series and parallel simplification | Easy]
\begin{enumerate}[label=(\alph*),topsep=2pt,itemsep=1pt,partopsep=0pt,parsep=0pt]
    \item {3 $\Omega$}
    \item {36 $\Omega$}
    \item {12 $\Omega$}
    \item {6 $\Omega$}
\end{enumerate}
\vspace{2mm}

\noindent\textbf{Q16. In a series circuit, the voltage division rule states that the voltage across a resistor is proportional to:} \hfill [Topic: voltage division rule | Easy]
\begin{enumerate}[label=(\alph*),topsep=2pt,itemsep=1pt,partopsep=0pt,parsep=0pt]
    \item {The source current only}
    \item {The reciprocal of its resistance}
    \item {The total current squared}
    \item {Its own resistance}
\end{enumerate}
\vspace{2mm}

\noindent\textbf{Q17. In a parallel circuit, the current division rule shows that more current flows through the branch with:} \hfill [Topic: current division rule | Easy]
\begin{enumerate}[label=(\alph*),topsep=2pt,itemsep=1pt,partopsep=0pt,parsep=0pt]
    \item {Higher resistance}
    \item {Lower resistance}
    \item {Longer wire length}
    \item {Larger inductance}
\end{enumerate}
\vspace{2mm}

\noindent\textbf{Q18. A source whose output depends on a voltage or current elsewhere in the circuit is called a:} \hfill [Topic: introduction to dependent and independent sources | Easy]
\begin{enumerate}[label=(\alph*),topsep=2pt,itemsep=1pt,partopsep=0pt,parsep=0pt]
    \item {Ideal source}
    \item {Short circuit}
    \item {Dependent source}
    \item {Independent source}
\end{enumerate}
\vspace{2mm}

\noindent\textbf{Q19. According to the maximum power transfer theorem, maximum power is delivered to the load when the load resistance equals the:} \hfill [Topic: maximum power transfer theorem | Easy]
\begin{enumerate}[label=(\alph*),topsep=2pt,itemsep=1pt,partopsep=0pt,parsep=0pt]
    \item {Twice the source resistance}
    \item {Total circuit current}
    \item {Source voltage}
    \item {Source (Thevenin) resistance}
\end{enumerate}
\vspace{2mm}

\noindent\textbf{Q20. The superposition theorem is applicable only to circuits that are:} \hfill [Topic: network theorems- superposition theorem | Easy]
\begin{enumerate}[label=(\alph*),topsep=2pt,itemsep=1pt,partopsep=0pt,parsep=0pt]
    \item {Linear}
    \item {Purely capacitive}
    \item {Nonlinear}
    \item {Open circuited}
\end{enumerate}
\vspace{2mm}

\noindent\textbf{Q21. A resistor carries a current of 0.5 A when connected across a 10 V supply. If the voltage is increased to 25 V , what current flows through the same resistor?} \hfill [Topic: ohm's law | Medium]
\begin{enumerate}[label=(\alph*),topsep=2pt,itemsep=1pt,partopsep=0pt,parsep=0pt]
    \item {0.5 A}
    \item {0.2 A}
    \item {2.5 A}
    \item {1.25 A}
\end{enumerate}
\vspace{2mm}

\noindent\textbf{Q22. Three resistors of 6 $\Omega$ , 3 $\Omega$ , and 2 $\Omega$ are connected in parallel. What is the equivalent resistance?} \hfill [Topic: resistance | Medium]
\begin{enumerate}[label=(\alph*),topsep=2pt,itemsep=1pt,partopsep=0pt,parsep=0pt]
    \item {3.67 $\Omega$}
    \item {1 $\Omega$}
    \item {11 $\Omega$}
    \item {0.5 $\Omega$}
\end{enumerate}
\vspace{2mm}

\noindent\textbf{Q23. A 100 $\Omega$ resistor dissipates 4 W of power. What is the current flowing through it?} \hfill [Topic: power and energy concepts | Medium]
\begin{enumerate}[label=(\alph*),topsep=2pt,itemsep=1pt,partopsep=0pt,parsep=0pt]
    \item {0.04 A}
    \item {0.2 A}
    \item {0.4 A}
    \item {2 A}
\end{enumerate}
\vspace{2mm}

\noindent\textbf{Q24. Two capacitors of 4 $\mu$ F and 12 $\mu$ F are connected in series. What is the equivalent capacitance?} \hfill [Topic: capacitance | Medium]
\begin{enumerate}[label=(\alph*),topsep=2pt,itemsep=1pt,partopsep=0pt,parsep=0pt]
    \item {16 $\mu$ F}
    \item {3 $\mu$ F}
    \item {8 $\mu$ F}
    \item {48 $\mu$ F}
\end{enumerate}
\vspace{2mm}

\noindent\textbf{Q25. The voltage across an inductor of 0.5 H when the current changes at a rate of 4 A / s is:} \hfill [Topic: inductance | Medium]
\begin{enumerate}[label=(\alph*),topsep=2pt,itemsep=1pt,partopsep=0pt,parsep=0pt]
    \item {2 V}
    \item {8 V}
    \item {4.5 V}
    \item {0.125 V}
\end{enumerate}
\vspace{2mm}

\noindent\textbf{Q26. In a series circuit, a 12 V source is connected across resistors of 3 $\Omega$ and 9 $\Omega$ . What is the voltage across the 9 $\Omega$ resistor?} \hfill [Topic: voltage division rule | Medium]
\begin{enumerate}[label=(\alph*),topsep=2pt,itemsep=1pt,partopsep=0pt,parsep=0pt]
    \item {9 V}
    \item {6 V}
    \item {3 V}
    \item {4 V}
\end{enumerate}
\vspace{2mm}

\noindent\textbf{Q27. A total current of 6 A enters two parallel resistors of 2 $\Omega$ and 4 $\Omega$ . What current flows through the 4 $\Omega$ resistor?} \hfill [Topic: current division rule | Medium]
\begin{enumerate}[label=(\alph*),topsep=2pt,itemsep=1pt,partopsep=0pt,parsep=0pt]
    \item {4 A}
    \item {3 A}
    \item {1.5 A}
    \item {2 A}
\end{enumerate}
\vspace{2mm}

\noindent\textbf{Q28. At a node, three currents enter: 5 A , 3 A , and 2 A , while one current leaves as 4 A . What is the magnitude of the second leaving current?} \hfill [Topic: kirchhoff's laws | Medium]
\begin{enumerate}[label=(\alph*),topsep=2pt,itemsep=1pt,partopsep=0pt,parsep=0pt]
    \item {2 A}
    \item {4 A}
    \item {10 A}
    \item {6 A}
\end{enumerate}
\vspace{2mm}

\noindent\textbf{Q29. In a closed loop, the sum of EMF sources is 20 V . If two resistors drop 8 V and 7 V , what is the voltage drop across the third resistor?} \hfill [Topic: kirchhoff's laws | Medium]
\begin{enumerate}[label=(\alph*),topsep=2pt,itemsep=1pt,partopsep=0pt,parsep=0pt]
    \item {15 V}
    \item {5 V}
    \item {12 V}
    \item {35 V}
\end{enumerate}
\vspace{2mm}

\noindent\textbf{Q30. A 6 $\Omega$ resistor is in series with a parallel combination of two 4 $\Omega$ resistors. What is the total resistance?} \hfill [Topic: intuitive method of circuit analysis- series and parallel simplification | Medium]
\begin{enumerate}[label=(\alph*),topsep=2pt,itemsep=1pt,partopsep=0pt,parsep=0pt]
    \item {10 $\Omega$}
    \item {8 $\Omega$}
    \item {14 $\Omega$}
    \item {2 $\Omega$}
\end{enumerate}
\vspace{2mm}

\noindent\textbf{Q31. Three equal resistors of 9 $\Omega$ are connected in delta. What is the resistance of each arm of the equivalent star network?} \hfill [Topic: star--delta transformation | Medium]
\begin{enumerate}[label=(\alph*),topsep=2pt,itemsep=1pt,partopsep=0pt,parsep=0pt]
    \item {3 $\Omega$}
    \item {27 $\Omega$}
    \item {9 $\Omega$}
    \item {6 $\Omega$}
\end{enumerate}
\vspace{2mm}

\noindent\textbf{Q32. Which of the following best describes a dependent source?} \hfill [Topic: introduction to dependent and independent sources | Medium]
\begin{enumerate}[label=(\alph*),topsep=2pt,itemsep=1pt,partopsep=0pt,parsep=0pt]
    \item {Its value is fixed regardless of circuit conditions}
    \item {It can only absorb power, never deliver it}
    \item {It supplies only alternating current}
    \item {Its value is controlled by a voltage or current elsewhere in the circuit}
\end{enumerate}
\vspace{2mm}

\noindent\textbf{Q33. For a planar circuit with 4 branches and 3 nodes, how many independent mesh equations are required?} \hfill [Topic: mesh and nodal analysis | Medium]
\begin{enumerate}[label=(\alph*),topsep=2pt,itemsep=1pt,partopsep=0pt,parsep=0pt]
    \item {1}
    \item {3}
    \item {4}
    \item {2}
\end{enumerate}
\vspace{2mm}

\noindent\textbf{Q34. In nodal analysis, how many node equations are needed for a circuit having 5 nodes (including the reference node)?} \hfill [Topic: mesh and nodal analysis | Medium]
\begin{enumerate}[label=(\alph*),topsep=2pt,itemsep=1pt,partopsep=0pt,parsep=0pt]
    \item {6}
    \item {4}
    \item {3}
    \item {5}
\end{enumerate}
\vspace{2mm}

\noindent\textbf{Q35. A network has an open-circuit voltage of 12 V and delivers a short-circuit current of 3 A at its terminals. What is the Thevenin resistance?} \hfill [Topic: Thevenin's theorem | Medium]
\begin{enumerate}[label=(\alph*),topsep=2pt,itemsep=1pt,partopsep=0pt,parsep=0pt]
    \item {9 $\Omega$}
    \item {0.25 $\Omega$}
    \item {36 $\Omega$}
    \item {4 $\Omega$}
\end{enumerate}
\vspace{2mm}

\noindent\textbf{Q36. A Thevenin equivalent has V t h = 20 V and R t h = 5 $\Omega$ . What is the Norton current of the equivalent circuit?} \hfill [Topic: Norton's theorem | Medium]
\begin{enumerate}[label=(\alph*),topsep=2pt,itemsep=1pt,partopsep=0pt,parsep=0pt]
    \item {100 A}
    \item {4 A}
    \item {0.25 A}
    \item {15 A}
\end{enumerate}
\vspace{2mm}

\noindent\textbf{Q37. A source has Thevenin resistance 8 $\Omega$ and open-circuit voltage 16 V . What is the maximum power delivered to a matched load?} \hfill [Topic: maximum power transfer theorem | Medium]
\begin{enumerate}[label=(\alph*),topsep=2pt,itemsep=1pt,partopsep=0pt,parsep=0pt]
    \item {32 W}
    \item {4 W}
    \item {16 W}
    \item {8 W}
\end{enumerate}
\vspace{2mm}

\noindent\textbf{Q38. For maximum power transfer to a resistive load in a DC circuit, the load resistance should be:} \hfill [Topic: maximum power transfer theorem | Medium]
\begin{enumerate}[label=(\alph*),topsep=2pt,itemsep=1pt,partopsep=0pt,parsep=0pt]
    \item {Zero}
    \item {Infinite}
    \item {Equal to the Thevenin resistance of the source}
    \item {Much larger than the source resistance}
\end{enumerate}
\vspace{2mm}

\noindent\textbf{Q39. While applying the superposition theorem, an independent voltage source that is not being considered must be:} \hfill [Topic: network theorems- superposition theorem | Medium]
\begin{enumerate}[label=(\alph*),topsep=2pt,itemsep=1pt,partopsep=0pt,parsep=0pt]
    \item {Doubled in value}
    \item {Left unchanged}
    \item {Replaced by an open circuit}
    \item {Replaced by a short circuit}
\end{enumerate}
\vspace{2mm}

\noindent\textbf{Q40. A charge of 120 C flows through a conductor in 2 minutes. What is the average current?} \hfill [Topic: current | Medium]
\begin{enumerate}[label=(\alph*),topsep=2pt,itemsep=1pt,partopsep=0pt,parsep=0pt]
    \item {2 A}
    \item {1 A}
    \item {60 A}
    \item {240 A}
\end{enumerate}
\vspace{2mm}

\noindent\textbf{Q41. A balanced delta network has three equal resistors of R $\Delta$ = 18 $\Omega$ each. If this is converted to an equivalent star network, what is the value of each star resistor?} \hfill [Topic: star--delta transformation | Hard]
\begin{enumerate}[label=(\alph*),topsep=2pt,itemsep=1pt,partopsep=0pt,parsep=0pt]
    \item {54 $\Omega$}
    \item {6 $\Omega$}
    \item {9 $\Omega$}
    \item {3 $\Omega$}
\end{enumerate}
\vspace{2mm}

\noindent\textbf{Q42. A network has an open-circuit voltage of 24 V and delivers a short-circuit current of 4 A at its terminals. A load resistor R L = 8 $\Omega$ is connected. What current flows through R L ?} \hfill [Topic: Thevenin's theorem | Hard]
\begin{enumerate}[label=(\alph*),topsep=2pt,itemsep=1pt,partopsep=0pt,parsep=0pt]
    \item {1 A}
    \item {3 A}
    \item {2 A}
    \item {1.71 A}
\end{enumerate}
\vspace{2mm}

\noindent\textbf{Q43. A source with Thevenin voltage 20 V and internal resistance 5 $\Omega$ drives a variable load. What is the maximum power that can be delivered to the load?} \hfill [Topic: maximum power transfer theorem | Hard]
\begin{enumerate}[label=(\alph*),topsep=2pt,itemsep=1pt,partopsep=0pt,parsep=0pt]
    \item {10 W}
    \item {40 W}
    \item {80 W}
    \item {20 W}
\end{enumerate}
\vspace{2mm}

\noindent\textbf{Q44. In a linear circuit, a 10 V source alone produces 2 A in a branch, and a 5 A source alone produces - 1 A in the same branch. With both sources active, what is the branch current?} \hfill [Topic: network theorems- superposition theorem | Hard]
\begin{enumerate}[label=(\alph*),topsep=2pt,itemsep=1pt,partopsep=0pt,parsep=0pt]
    \item {3 A}
    \item {1 A}
    \item {- 3 A}
    \item {2 A}
\end{enumerate}
\vspace{2mm}

\noindent\textbf{Q45. A two-node circuit (plus reference) has node equations: 3 V 1 - V 2 = 10 and - V 1 + 2 V 2 = 5 (conductances in S, currents in A). What is V 1 ?} \hfill [Topic: mesh and nodal analysis | Hard]
\begin{enumerate}[label=(\alph*),topsep=2pt,itemsep=1pt,partopsep=0pt,parsep=0pt]
    \item {2.5 V}
    \item {4 V}
    \item {5 V}
    \item {3 V}
\end{enumerate}
\vspace{2mm}

\noindent\textbf{Q46. A total current of 12 A enters two parallel resistors R 1 = 4 $\Omega$ and R 2 = 12 $\Omega$ . What current flows through R 1 ?} \hfill [Topic: current division rule | Hard]
\begin{enumerate}[label=(\alph*),topsep=2pt,itemsep=1pt,partopsep=0pt,parsep=0pt]
    \item {3 A}
    \item {9 A}
    \item {6 A}
    \item {8 A}
\end{enumerate}
\vspace{2mm}

\noindent\textbf{Q47. Three resistors 2 $\Omega$ , 3 $\Omega$ , and 5 $\Omega$ are in series across a 20 V source. What voltage appears across the 3 $\Omega$ resistor?} \hfill [Topic: voltage division rule | Hard]
\begin{enumerate}[label=(\alph*),topsep=2pt,itemsep=1pt,partopsep=0pt,parsep=0pt]
    \item {3 V}
    \item {4 V}
    \item {10 V}
    \item {6 V}
\end{enumerate}
\vspace{2mm}

\noindent\textbf{Q48. At a node, three currents meet: I 1 = 5 A entering, I 2 = 3 A leaving, and I 3 unknown. If KCL holds and I 3 leaves the node, what is I 3 ?} \hfill [Topic: kirchhoff's laws | Hard]
\begin{enumerate}[label=(\alph*),topsep=2pt,itemsep=1pt,partopsep=0pt,parsep=0pt]
    \item {2 A}
    \item {3 A}
    \item {- 2 A}
    \item {8 A}
\end{enumerate}
\vspace{2mm}

\noindent\textbf{Q49. A Norton equivalent has I N = 6 A and R N = 4 $\Omega$ . A load of 4 $\Omega$ is connected. What power is dissipated in the load?} \hfill [Topic: Norton's theorem | Hard]
\begin{enumerate}[label=(\alph*),topsep=2pt,itemsep=1pt,partopsep=0pt,parsep=0pt]
    \item {18 W}
    \item {144 W}
    \item {36 W}
    \item {72 W}
\end{enumerate}
\vspace{2mm}

\noindent\textbf{Q50. Two capacitors C 1 = 6 $\mu$ F and C 2 = 3 $\mu$ F are connected in series across 30 V . What is the charge stored on each capacitor?} \hfill [Topic: capacitance | Hard]
\begin{enumerate}[label=(\alph*),topsep=2pt,itemsep=1pt,partopsep=0pt,parsep=0pt]
    \item {180 $\mu$ C}
    \item {60 $\mu$ C}
    \item {90 $\mu$ C}
    \item {30 $\mu$ C}
\end{enumerate}
\vspace{2mm}

\noindent\textbf{Q51. The current through a 0.5 H inductor changes from 2 A to 6 A in 0.1 s . What is the magnitude of the induced EMF?} \hfill [Topic: inductance | Hard]
\begin{enumerate}[label=(\alph*),topsep=2pt,itemsep=1pt,partopsep=0pt,parsep=0pt]
    \item {5 V}
    \item {40 V}
    \item {20 V}
    \item {10 V}
\end{enumerate}
\vspace{2mm}

\noindent\textbf{Q52. A 100 W device operates for 5 hours daily. What is the energy consumed in one week in kWh?} \hfill [Topic: power and energy concepts | Hard]
\begin{enumerate}[label=(\alph*),topsep=2pt,itemsep=1pt,partopsep=0pt,parsep=0pt]
    \item {0.5 kWh}
    \item {35 kWh}
    \item {3.5 kWh}
    \item {7 kWh}
\end{enumerate}
\vspace{2mm}

\noindent\textbf{Q53. A wire of resistance R is stretched uniformly to twice its original length (volume constant). What is its new resistance?} \hfill [Topic: resistance | Hard]
\begin{enumerate}[label=(\alph*),topsep=2pt,itemsep=1pt,partopsep=0pt,parsep=0pt]
    \item {R /2}
    \item {R /4}
    \item {2 R}
    \item {4 R}
\end{enumerate}
\vspace{2mm}

\noindent\textbf{Q54. A resistor dissipates 50 W when carrying 5 A . What voltage is across it?} \hfill [Topic: ohm's law | Hard]
\begin{enumerate}[label=(\alph*),topsep=2pt,itemsep=1pt,partopsep=0pt,parsep=0pt]
    \item {250 V}
    \item {10 V}
    \item {55 V}
    \item {2 V}
\end{enumerate}
\vspace{2mm}

\noindent\textbf{Q55. Two 6 $\Omega$ resistors in parallel are placed in series with a 4 $\Omega$ resistor, and this combination is in parallel with a 7 $\Omega$ resistor. What is the total resistance?} \hfill [Topic: intuitive method of circuit analysis- series and parallel simplification | Hard]
\begin{enumerate}[label=(\alph*),topsep=2pt,itemsep=1pt,partopsep=0pt,parsep=0pt]
    \item {3.5 $\Omega$}
    \item {7 $\Omega$}
    \item {10 $\Omega$}
    \item {2.33 $\Omega$}
\end{enumerate}
\vspace{2mm}

\noindent\textbf{Q56. A circuit contains a current-controlled voltage source V s = 4 I x where I x is the current through a 2 $\Omega$ resistor carrying 3 A . What is the source voltage?} \hfill [Topic: introduction to dependent and independent sources | Hard]
\begin{enumerate}[label=(\alph*),topsep=2pt,itemsep=1pt,partopsep=0pt,parsep=0pt]
    \item {12 V}
    \item {6 V}
    \item {8 V}
    \item {24 V}
\end{enumerate}
\vspace{2mm}

\noindent\textbf{Q57. A circuit has 4 nodes and 6 branches. How many independent KVL (loop) equations are needed for complete analysis?} \hfill [Topic: basic method of circuit analysis | Hard]
\begin{enumerate}[label=(\alph*),topsep=2pt,itemsep=1pt,partopsep=0pt,parsep=0pt]
    \item {2}
    \item {3}
    \item {6}
    \item {4}
\end{enumerate}
\vspace{2mm}

\noindent\textbf{Q58. Moving a charge of 4 C between two points requires 48 J of work. What is the potential difference between the points?} \hfill [Topic: voltage | Hard]
\begin{enumerate}[label=(\alph*),topsep=2pt,itemsep=1pt,partopsep=0pt,parsep=0pt]
    \item {192 V}
    \item {44 V}
    \item {0.083 V}
    \item {12 V}
\end{enumerate}
\vspace{2mm}

\noindent\textbf{Q59. A charge of 600 C passes through a conductor in 5 minutes. What is the average current?} \hfill [Topic: current | Hard]
\begin{enumerate}[label=(\alph*),topsep=2pt,itemsep=1pt,partopsep=0pt,parsep=0pt]
    \item {10 A}
    \item {2 A}
    \item {0.5 A}
    \item {120 A}
\end{enumerate}
\vspace{2mm}

\noindent\textbf{Q60. A two-mesh circuit yields: 10 = 6 I 1 - 2 I 2 and 0 = - 2 I 1 + 5 I 2 . What is the mesh current I 1 ?} \hfill [Topic: mesh and nodal analysis | Hard]
\begin{enumerate}[label=(\alph*),topsep=2pt,itemsep=1pt,partopsep=0pt,parsep=0pt]
    \item {1.67 A}
    \item {0.77 A}
    \item {1.92 A}
    \item {2.5 A}
\end{enumerate}
\vspace{2mm}


\newpage
\section{Section B: Comprehensive Analytical Problems with Step-by-Step Solutions}

\begin{subjectiveq}{1}{1}
In the circuit shown below, determine the branch currents $I_1$, $I_2$, and $I_3$ using Mesh Current Analysis. Also determine the Thevenin equivalent circuit across terminals $A-B$ connected across the $10\,\Omega$ load resistor, and find the maximum power transferable to $R_L$.\\[2mm]Circuit Parameters: Independent voltage sources $V_1 = 36\text{ V}$, $V_2 = 12\text{ V}$; Resistors $R_1 = 6\,\Omega$, $R_2 = 3\,\Omega$, $R_3 = 6\,\Omega$, Load $R_L = 10\,\Omega$.
\end{subjectiveq}

\begin{solbox}
\textbf{Step 1: Mesh Formulation and Matrix Setup}\\Define clockwise mesh currents $I_a$ in Loop 1 (left loop) and $I_b$ in Loop 2 (right loop):\begin{align*}\text{Loop 1: } & (R_1 + R_2)I_a - R_2 I_b = V_1 \implies 9 I_a - 3 I_b = 36 \\\text{Loop 2: } & -R_2 I_a + (R_2 + R_3)I_b = -V_2 \implies -3 I_a + 9 I_b = -12\end{align*}\textbf{Step 2: Solving for Mesh Currents via Cramer's Rule}\\The determinant of the impedance matrix is $\Delta = (9)(9) - (-3)(-3) = 81 - 9 = 72\,\Omega^2$.\begin{align*}I_a &= \frac{\begin{vmatrix} 36 & -3 \\ -12 & 9 \end{vmatrix}}{72} = \frac{(36)(9) - (-12)(-3)}{72} = \frac{324 - 36}{72} = \frac{288}{72} = 4.0\text{ A} \\I_b &= \frac{\begin{vmatrix} 9 & 36 \\ -3 & -12 \end{vmatrix}}{72} = \frac{(9)(-12) - (-3)(36)}{72} = \frac{-108 + 108}{72} = 0.0\text{ A}\end{align*}Branch currents are therefore: $I_{R1} = I_a = 4.0\text{ A}$, $I_{R2} = I_a - I_b = 4.0\text{ A}$, and $I_{R3} = I_b = 0.0\text{ A}$.\\[2mm]\textbf{Step 3: Thevenin Equivalent ($V_{th}, R_{th}$ across terminals of $R_3$)}\\Open-circuiting $R_3$, the open circuit voltage across terminals $A-B$ is governed by the voltage divider of $V_1$ across $R_2$ minus $V_2$:\[ V_{th} = V_1 \frac{R_2}{R_1 + R_2} - V_2 = 36 \times \frac{3}{9} - 12 = 12 - 12 = 0\text{ V} \]Deactivating independent sources (short-circuiting $V_1$ and $V_2$):\[ R_{th} = (R_1 \parallel R_2) = \frac{6 \times 3}{6 + 3} = 2.0\,\Omega \]For Maximum Power Transfer to load $R_L$, we require $R_L = R_{th} = 2.0\,\Omega$.\[ P_{max} = \frac{V_{th}^2}{4 R_{th}} \]
\end{solbox}

\begin{examinertip}
\textbf{Grading Breakdown [Total: 10 Marks]}:\begin{itemize}[leftmargin=1.5em]\item Correct KVL mesh equations with clear sign conventions: \textbf{3 Marks}\item Matrix inversion / determinant computation of branch currents: \textbf{3 Marks}\item Thevenin resistance $R_{th}$ calculation with deactivated sources: \textbf{2 Marks}\item Maximum power theorem condition and final calculation: \textbf{2 Marks}\end{itemize}
\end{examinertip}
\vspace{4mm}
